% =========================================================
% Slides Beamer Metropolis 16:9 — Chapter 4 (Non-stationarity)
% Time Series Analysis — Aymen Ben Brik (aymen.benbrik@esprit.tn)
%
% Migrated from _archives/source-slides/Non stationnarité.tex (1059 L)
% sections 1-3 (non-stat + operators + ARIMA + SARIMA) and from
% Stationnarité.tex section "Tests de stationnarité" (lines 475-765).
% =========================================================
\documentclass[aspectratio=169]{beamer}

% =========================================================
% Beamer preamble — Time Series Analysis module
% Theme : Metropolis, aspect ratio 16:9
% Author : Aymen Ben Brik (aymen.benbrik@esprit.tn)
% Esprit School of Business
%
% Include from each chapterX/slides.tex via :
%   % =========================================================
% Beamer preamble — Time Series Analysis module
% Theme : Metropolis, aspect ratio 16:9
% Author : Aymen Ben Brik (aymen.benbrik@esprit.tn)
% Esprit School of Business
%
% Include from each chapterX/slides.tex via :
%   % =========================================================
% Beamer preamble — Time Series Analysis module
% Theme : Metropolis, aspect ratio 16:9
% Author : Aymen Ben Brik (aymen.benbrik@esprit.tn)
% Esprit School of Business
%
% Include from each chapterX/slides.tex via :
%   \input{../preamble/slides-preamble.tex}
% AFTER the line \documentclass[aspectratio=169]{beamer}.
% =========================================================

\usepackage[utf8]{inputenc}
\usepackage[T1]{fontenc}
\usepackage[english]{babel}
\usepackage{lmodern}
\usepackage{amsmath, amssymb, mathtools, bm}
\usepackage{graphicx}
\usepackage{booktabs}
\usepackage{array}
\usepackage{multirow}
\usepackage{colortbl}
\usepackage{xcolor}
\usepackage{tikz}
\usetikzlibrary{positioning, arrows.meta, calc, shapes, fit, backgrounds, matrix, decorations.pathreplacing}
\usepackage{listings}
\usepackage[most]{tcolorbox}

% --- Theme ---
\usetheme{metropolis}
\setbeamertemplate{navigation symbols}{}
\metroset{block=fill}

% --- Esprit colour palette ---
\definecolor{espritBlue}{RGB}{30, 60, 110}
\definecolor{espritAccent}{RGB}{200, 80, 30}
\setbeamercolor{progress bar}{fg=espritAccent}
\setbeamercolor{title separator}{fg=espritAccent}
\setbeamercolor{frametitle}{bg=espritBlue}

% --- Code listings (Python + R, with UTF-8 literate) ---
\definecolor{codebg}{RGB}{248,248,248}
\definecolor{codeframe}{RGB}{220,220,220}
\definecolor{keyword}{RGB}{0,82,155}
\definecolor{comment}{RGB}{88,110,117}
\definecolor{string}{RGB}{133,153,0}

\lstdefinestyle{pythonstyle}{
  language=Python,
  basicstyle=\ttfamily\scriptsize,
  keywordstyle=\color{keyword}\bfseries,
  commentstyle=\color{comment}\itshape,
  stringstyle=\color{string},
  backgroundcolor=\color{codebg},
  frame=single,
  rulecolor=\color{codeframe},
  frameround=tttt,
  breaklines=true,
  showstringspaces=false,
  tabsize=2,
  literate=
    {é}{{\'e}}1 {è}{{\`e}}1 {ê}{{\^e}}1 {à}{{\`a}}1 {â}{{\^a}}1
    {ç}{{\c{c}}}1 {²}{{$^2$}}1 {°}{{${}^\circ$}}1
    {α}{{$\alpha$}}1 {β}{{$\beta$}}1 {γ}{{$\gamma$}}1
    {δ}{{$\delta$}}1 {θ}{{$\theta$}}1 {λ}{{$\lambda$}}1
    {μ}{{$\mu$}}1 {σ}{{$\sigma$}}1 {ε}{{$\varepsilon$}}1
    {ρ}{{$\rho$}}1 {τ}{{$\tau$}}1 {φ}{{$\phi$}}1
    {≤}{{$\leq$}}1 {≥}{{$\geq$}}1 {≈}{{$\approx$}}1
}
\lstdefinestyle{Rstyle}{
  language=R,
  basicstyle=\ttfamily\scriptsize,
  keywordstyle=\color{keyword}\bfseries,
  commentstyle=\color{comment}\itshape,
  stringstyle=\color{string},
  backgroundcolor=\color{codebg},
  frame=single,
  rulecolor=\color{codeframe},
  frameround=tttt,
  breaklines=true,
  showstringspaces=false,
  tabsize=2
}
\lstset{style=pythonstyle}

% --- Common macros ---
\input{../preamble/macros.tex}

% --- tcolorbox boxes for slides ---
\definecolor{boxTh}{RGB}{120, 30, 30}
\definecolor{boxProp}{RGB}{30, 100, 60}
\definecolor{boxDef}{RGB}{30, 60, 110}
\definecolor{boxEx}{RGB}{180, 120, 0}

\newtcolorbox{infobox}[1][]{enhanced, colback=espritBlue!5, colframe=espritBlue,
  fonttitle=\bfseries\small,
  title={\ifstrempty{#1}{}{#1}},
  boxrule=0.4pt, arc=2pt, left=4pt, right=4pt, top=2pt, bottom=2pt}
\newtcolorbox{warnbox}[1][]{enhanced, colback=espritAccent!5, colframe=espritAccent,
  fonttitle=\bfseries\small,
  title={\ifstrempty{#1}{}{#1}},
  boxrule=0.4pt, arc=2pt, left=4pt, right=4pt, top=2pt, bottom=2pt}
\newtcolorbox{thbox}[1][]{enhanced, colback=boxTh!5, colframe=boxTh,
  fonttitle=\bfseries\small,
  title={Theorem\ifstrempty{#1}{}{ -- #1}},
  boxrule=0.5pt, arc=2pt, left=4pt, right=4pt, top=2pt, bottom=2pt}
\newtcolorbox{propbox}[1][]{enhanced, colback=boxProp!5, colframe=boxProp,
  fonttitle=\bfseries\small,
  title={Proposition\ifstrempty{#1}{}{ -- #1}},
  boxrule=0.5pt, arc=2pt, left=4pt, right=4pt, top=2pt, bottom=2pt}
\newtcolorbox{defbox}[1][]{enhanced, colback=boxDef!5, colframe=boxDef,
  fonttitle=\bfseries\small,
  title={Definition\ifstrempty{#1}{}{ -- #1}},
  boxrule=0.5pt, arc=2pt, left=4pt, right=4pt, top=2pt, bottom=2pt}
\newtcolorbox{exbox}[1][]{enhanced, colback=boxEx!5, colframe=boxEx,
  fonttitle=\bfseries\small,
  title={Example\ifstrempty{#1}{}{ -- #1}},
  boxrule=0.5pt, arc=2pt, left=4pt, right=4pt, top=2pt, bottom=2pt}

% AFTER the line \documentclass[aspectratio=169]{beamer}.
% =========================================================

\usepackage[utf8]{inputenc}
\usepackage[T1]{fontenc}
\usepackage[english]{babel}
\usepackage{lmodern}
\usepackage{amsmath, amssymb, mathtools, bm}
\usepackage{graphicx}
\usepackage{booktabs}
\usepackage{array}
\usepackage{multirow}
\usepackage{colortbl}
\usepackage{xcolor}
\usepackage{tikz}
\usetikzlibrary{positioning, arrows.meta, calc, shapes, fit, backgrounds, matrix, decorations.pathreplacing}
\usepackage{listings}
\usepackage[most]{tcolorbox}

% --- Theme ---
\usetheme{metropolis}
\setbeamertemplate{navigation symbols}{}
\metroset{block=fill}

% --- Esprit colour palette ---
\definecolor{espritBlue}{RGB}{30, 60, 110}
\definecolor{espritAccent}{RGB}{200, 80, 30}
\setbeamercolor{progress bar}{fg=espritAccent}
\setbeamercolor{title separator}{fg=espritAccent}
\setbeamercolor{frametitle}{bg=espritBlue}

% --- Code listings (Python + R, with UTF-8 literate) ---
\definecolor{codebg}{RGB}{248,248,248}
\definecolor{codeframe}{RGB}{220,220,220}
\definecolor{keyword}{RGB}{0,82,155}
\definecolor{comment}{RGB}{88,110,117}
\definecolor{string}{RGB}{133,153,0}

\lstdefinestyle{pythonstyle}{
  language=Python,
  basicstyle=\ttfamily\scriptsize,
  keywordstyle=\color{keyword}\bfseries,
  commentstyle=\color{comment}\itshape,
  stringstyle=\color{string},
  backgroundcolor=\color{codebg},
  frame=single,
  rulecolor=\color{codeframe},
  frameround=tttt,
  breaklines=true,
  showstringspaces=false,
  tabsize=2,
  literate=
    {é}{{\'e}}1 {è}{{\`e}}1 {ê}{{\^e}}1 {à}{{\`a}}1 {â}{{\^a}}1
    {ç}{{\c{c}}}1 {²}{{$^2$}}1 {°}{{${}^\circ$}}1
    {α}{{$\alpha$}}1 {β}{{$\beta$}}1 {γ}{{$\gamma$}}1
    {δ}{{$\delta$}}1 {θ}{{$\theta$}}1 {λ}{{$\lambda$}}1
    {μ}{{$\mu$}}1 {σ}{{$\sigma$}}1 {ε}{{$\varepsilon$}}1
    {ρ}{{$\rho$}}1 {τ}{{$\tau$}}1 {φ}{{$\phi$}}1
    {≤}{{$\leq$}}1 {≥}{{$\geq$}}1 {≈}{{$\approx$}}1
}
\lstdefinestyle{Rstyle}{
  language=R,
  basicstyle=\ttfamily\scriptsize,
  keywordstyle=\color{keyword}\bfseries,
  commentstyle=\color{comment}\itshape,
  stringstyle=\color{string},
  backgroundcolor=\color{codebg},
  frame=single,
  rulecolor=\color{codeframe},
  frameround=tttt,
  breaklines=true,
  showstringspaces=false,
  tabsize=2
}
\lstset{style=pythonstyle}

% --- Common macros ---
% =========================================================
% Common math macros — Time Series Analysis module
% Author : Aymen Ben Brik (aymen.benbrik@esprit.tn)
% Esprit School of Business
% =========================================================

% --- Standard sets ---
\providecommand{\R}{\mathbb{R}}
\providecommand{\N}{\mathbb{N}}
\providecommand{\Z}{\mathbb{Z}}
\providecommand{\Q}{\mathbb{Q}}
\providecommand{\C}{\mathbb{C}}

% --- Probability / statistics ---
\providecommand{\E}{\mathbb{E}}
\providecommand{\Prob}{\mathbb{P}}
\providecommand{\Var}{\mathrm{Var}}
\providecommand{\Cov}{\mathrm{Cov}}
\providecommand{\Corr}{\mathrm{Corr}}
\providecommand{\indep}{\perp\!\!\!\perp}
\providecommand{\iid}{\overset{\text{i.i.d.}}{\sim}}
\providecommand{\given}{\,\big\vert\,}

% --- Estimators / hat ---
\providecommand{\hatm}{\widehat{m}}
\providecommand{\hatsig}{\widehat{\sigma}}
\providecommand{\hatrho}{\widehat{\rho}}
\providecommand{\hatphi}{\widehat{\phi}}
\providecommand{\hattheta}{\widehat{\theta}}

% --- Time series notation ---
\providecommand{\Lop}{L}                      % lag operator
\providecommand{\Diff}{\Delta}                % difference operator
\providecommand{\AR}[1]{\mathrm{AR}(#1)}
\providecommand{\MAm}[1]{\mathrm{MA}(#1)}
\providecommand{\ARMA}[2]{\mathrm{ARMA}(#1,#2)}
\providecommand{\ARIMA}[3]{\mathrm{ARIMA}(#1,#2,#3)}
\providecommand{\SARIMA}[6]{\mathrm{SARIMA}(#1,#2,#3)(#4,#5,#6)}
\providecommand{\ARCH}[1]{\mathrm{ARCH}(#1)}
\providecommand{\GARCH}[2]{\mathrm{GARCH}(#1,#2)}
\providecommand{\WN}{\mathrm{WN}}              % white noise
\providecommand{\MSE}{\mathrm{MSE}}
\providecommand{\MAE}{\mathrm{MAE}}
\providecommand{\RMSE}{\mathrm{RMSE}}
\providecommand{\AIC}{\mathrm{AIC}}
\providecommand{\BIC}{\mathrm{BIC}}
\providecommand{\ACF}{\mathrm{ACF}}
\providecommand{\PACF}{\mathrm{PACF}}

% --- Linear algebra ---
\providecommand{\transp}[1]{#1^{\!\top}}
\providecommand{\norm}[1]{\left\lVert #1 \right\rVert}
\providecommand{\abs}[1]{\left\lvert #1 \right\rvert}
\providecommand{\inner}[2]{\langle #1, #2 \rangle}
\providecommand{\diag}{\mathrm{diag}}
\providecommand{\tr}{\mathrm{tr}}

% --- Bold vectors / matrices ---
\providecommand{\vx}{\bm{x}}
\providecommand{\vy}{\bm{y}}
\providecommand{\vz}{\bm{z}}
\providecommand{\vh}{\bm{h}}
\providecommand{\vu}{\bm{u}}
\providecommand{\vv}{\bm{v}}
\providecommand{\vw}{\bm{w}}
\providecommand{\vb}{\bm{b}}
\providecommand{\vW}{\bm{W}}
\providecommand{\vSigma}{\bm{\Sigma}}
\providecommand{\veps}{\bm{\varepsilon}}

% --- Author identity (reused on titlepages) ---
\providecommand{\auteurNom}{Aymen Ben Brik}
\providecommand{\auteurEmail}{aymen.benbrik@esprit.tn}
\providecommand{\auteurInstitution}{Esprit School of Business}
\providecommand{\auteurRole}{Module head}
\providecommand{\moduleNom}{Time Series Analysis}
\providecommand{\anneeUniversitaire}{2025--2026}


% --- tcolorbox boxes for slides ---
\definecolor{boxTh}{RGB}{120, 30, 30}
\definecolor{boxProp}{RGB}{30, 100, 60}
\definecolor{boxDef}{RGB}{30, 60, 110}
\definecolor{boxEx}{RGB}{180, 120, 0}

\newtcolorbox{infobox}[1][]{enhanced, colback=espritBlue!5, colframe=espritBlue,
  fonttitle=\bfseries\small,
  title={\ifstrempty{#1}{}{#1}},
  boxrule=0.4pt, arc=2pt, left=4pt, right=4pt, top=2pt, bottom=2pt}
\newtcolorbox{warnbox}[1][]{enhanced, colback=espritAccent!5, colframe=espritAccent,
  fonttitle=\bfseries\small,
  title={\ifstrempty{#1}{}{#1}},
  boxrule=0.4pt, arc=2pt, left=4pt, right=4pt, top=2pt, bottom=2pt}
\newtcolorbox{thbox}[1][]{enhanced, colback=boxTh!5, colframe=boxTh,
  fonttitle=\bfseries\small,
  title={Theorem\ifstrempty{#1}{}{ -- #1}},
  boxrule=0.5pt, arc=2pt, left=4pt, right=4pt, top=2pt, bottom=2pt}
\newtcolorbox{propbox}[1][]{enhanced, colback=boxProp!5, colframe=boxProp,
  fonttitle=\bfseries\small,
  title={Proposition\ifstrempty{#1}{}{ -- #1}},
  boxrule=0.5pt, arc=2pt, left=4pt, right=4pt, top=2pt, bottom=2pt}
\newtcolorbox{defbox}[1][]{enhanced, colback=boxDef!5, colframe=boxDef,
  fonttitle=\bfseries\small,
  title={Definition\ifstrempty{#1}{}{ -- #1}},
  boxrule=0.5pt, arc=2pt, left=4pt, right=4pt, top=2pt, bottom=2pt}
\newtcolorbox{exbox}[1][]{enhanced, colback=boxEx!5, colframe=boxEx,
  fonttitle=\bfseries\small,
  title={Example\ifstrempty{#1}{}{ -- #1}},
  boxrule=0.5pt, arc=2pt, left=4pt, right=4pt, top=2pt, bottom=2pt}

% AFTER the line \documentclass[aspectratio=169]{beamer}.
% =========================================================

\usepackage[utf8]{inputenc}
\usepackage[T1]{fontenc}
\usepackage[english]{babel}
\usepackage{lmodern}
\usepackage{amsmath, amssymb, mathtools, bm}
\usepackage{graphicx}
\usepackage{booktabs}
\usepackage{array}
\usepackage{multirow}
\usepackage{colortbl}
\usepackage{xcolor}
\usepackage{tikz}
\usetikzlibrary{positioning, arrows.meta, calc, shapes, fit, backgrounds, matrix, decorations.pathreplacing}
\usepackage{listings}
\usepackage[most]{tcolorbox}

% --- Theme ---
\usetheme{metropolis}
\setbeamertemplate{navigation symbols}{}
\metroset{block=fill}

% --- Esprit colour palette ---
\definecolor{espritBlue}{RGB}{30, 60, 110}
\definecolor{espritAccent}{RGB}{200, 80, 30}
\setbeamercolor{progress bar}{fg=espritAccent}
\setbeamercolor{title separator}{fg=espritAccent}
\setbeamercolor{frametitle}{bg=espritBlue}

% --- Code listings (Python + R, with UTF-8 literate) ---
\definecolor{codebg}{RGB}{248,248,248}
\definecolor{codeframe}{RGB}{220,220,220}
\definecolor{keyword}{RGB}{0,82,155}
\definecolor{comment}{RGB}{88,110,117}
\definecolor{string}{RGB}{133,153,0}

\lstdefinestyle{pythonstyle}{
  language=Python,
  basicstyle=\ttfamily\scriptsize,
  keywordstyle=\color{keyword}\bfseries,
  commentstyle=\color{comment}\itshape,
  stringstyle=\color{string},
  backgroundcolor=\color{codebg},
  frame=single,
  rulecolor=\color{codeframe},
  frameround=tttt,
  breaklines=true,
  showstringspaces=false,
  tabsize=2,
  literate=
    {é}{{\'e}}1 {è}{{\`e}}1 {ê}{{\^e}}1 {à}{{\`a}}1 {â}{{\^a}}1
    {ç}{{\c{c}}}1 {²}{{$^2$}}1 {°}{{${}^\circ$}}1
    {α}{{$\alpha$}}1 {β}{{$\beta$}}1 {γ}{{$\gamma$}}1
    {δ}{{$\delta$}}1 {θ}{{$\theta$}}1 {λ}{{$\lambda$}}1
    {μ}{{$\mu$}}1 {σ}{{$\sigma$}}1 {ε}{{$\varepsilon$}}1
    {ρ}{{$\rho$}}1 {τ}{{$\tau$}}1 {φ}{{$\phi$}}1
    {≤}{{$\leq$}}1 {≥}{{$\geq$}}1 {≈}{{$\approx$}}1
}
\lstdefinestyle{Rstyle}{
  language=R,
  basicstyle=\ttfamily\scriptsize,
  keywordstyle=\color{keyword}\bfseries,
  commentstyle=\color{comment}\itshape,
  stringstyle=\color{string},
  backgroundcolor=\color{codebg},
  frame=single,
  rulecolor=\color{codeframe},
  frameround=tttt,
  breaklines=true,
  showstringspaces=false,
  tabsize=2
}
\lstset{style=pythonstyle}

% --- Common macros ---
% =========================================================
% Common math macros — Time Series Analysis module
% Author : Aymen Ben Brik (aymen.benbrik@esprit.tn)
% Esprit School of Business
% =========================================================

% --- Standard sets ---
\providecommand{\R}{\mathbb{R}}
\providecommand{\N}{\mathbb{N}}
\providecommand{\Z}{\mathbb{Z}}
\providecommand{\Q}{\mathbb{Q}}
\providecommand{\C}{\mathbb{C}}

% --- Probability / statistics ---
\providecommand{\E}{\mathbb{E}}
\providecommand{\Prob}{\mathbb{P}}
\providecommand{\Var}{\mathrm{Var}}
\providecommand{\Cov}{\mathrm{Cov}}
\providecommand{\Corr}{\mathrm{Corr}}
\providecommand{\indep}{\perp\!\!\!\perp}
\providecommand{\iid}{\overset{\text{i.i.d.}}{\sim}}
\providecommand{\given}{\,\big\vert\,}

% --- Estimators / hat ---
\providecommand{\hatm}{\widehat{m}}
\providecommand{\hatsig}{\widehat{\sigma}}
\providecommand{\hatrho}{\widehat{\rho}}
\providecommand{\hatphi}{\widehat{\phi}}
\providecommand{\hattheta}{\widehat{\theta}}

% --- Time series notation ---
\providecommand{\Lop}{L}                      % lag operator
\providecommand{\Diff}{\Delta}                % difference operator
\providecommand{\AR}[1]{\mathrm{AR}(#1)}
\providecommand{\MAm}[1]{\mathrm{MA}(#1)}
\providecommand{\ARMA}[2]{\mathrm{ARMA}(#1,#2)}
\providecommand{\ARIMA}[3]{\mathrm{ARIMA}(#1,#2,#3)}
\providecommand{\SARIMA}[6]{\mathrm{SARIMA}(#1,#2,#3)(#4,#5,#6)}
\providecommand{\ARCH}[1]{\mathrm{ARCH}(#1)}
\providecommand{\GARCH}[2]{\mathrm{GARCH}(#1,#2)}
\providecommand{\WN}{\mathrm{WN}}              % white noise
\providecommand{\MSE}{\mathrm{MSE}}
\providecommand{\MAE}{\mathrm{MAE}}
\providecommand{\RMSE}{\mathrm{RMSE}}
\providecommand{\AIC}{\mathrm{AIC}}
\providecommand{\BIC}{\mathrm{BIC}}
\providecommand{\ACF}{\mathrm{ACF}}
\providecommand{\PACF}{\mathrm{PACF}}

% --- Linear algebra ---
\providecommand{\transp}[1]{#1^{\!\top}}
\providecommand{\norm}[1]{\left\lVert #1 \right\rVert}
\providecommand{\abs}[1]{\left\lvert #1 \right\rvert}
\providecommand{\inner}[2]{\langle #1, #2 \rangle}
\providecommand{\diag}{\mathrm{diag}}
\providecommand{\tr}{\mathrm{tr}}

% --- Bold vectors / matrices ---
\providecommand{\vx}{\bm{x}}
\providecommand{\vy}{\bm{y}}
\providecommand{\vz}{\bm{z}}
\providecommand{\vh}{\bm{h}}
\providecommand{\vu}{\bm{u}}
\providecommand{\vv}{\bm{v}}
\providecommand{\vw}{\bm{w}}
\providecommand{\vb}{\bm{b}}
\providecommand{\vW}{\bm{W}}
\providecommand{\vSigma}{\bm{\Sigma}}
\providecommand{\veps}{\bm{\varepsilon}}

% --- Author identity (reused on titlepages) ---
\providecommand{\auteurNom}{Aymen Ben Brik}
\providecommand{\auteurEmail}{aymen.benbrik@esprit.tn}
\providecommand{\auteurInstitution}{Esprit School of Business}
\providecommand{\auteurRole}{Module head}
\providecommand{\moduleNom}{Time Series Analysis}
\providecommand{\anneeUniversitaire}{2025--2026}


% --- tcolorbox boxes for slides ---
\definecolor{boxTh}{RGB}{120, 30, 30}
\definecolor{boxProp}{RGB}{30, 100, 60}
\definecolor{boxDef}{RGB}{30, 60, 110}
\definecolor{boxEx}{RGB}{180, 120, 0}

\newtcolorbox{infobox}[1][]{enhanced, colback=espritBlue!5, colframe=espritBlue,
  fonttitle=\bfseries\small,
  title={\ifstrempty{#1}{}{#1}},
  boxrule=0.4pt, arc=2pt, left=4pt, right=4pt, top=2pt, bottom=2pt}
\newtcolorbox{warnbox}[1][]{enhanced, colback=espritAccent!5, colframe=espritAccent,
  fonttitle=\bfseries\small,
  title={\ifstrempty{#1}{}{#1}},
  boxrule=0.4pt, arc=2pt, left=4pt, right=4pt, top=2pt, bottom=2pt}
\newtcolorbox{thbox}[1][]{enhanced, colback=boxTh!5, colframe=boxTh,
  fonttitle=\bfseries\small,
  title={Theorem\ifstrempty{#1}{}{ -- #1}},
  boxrule=0.5pt, arc=2pt, left=4pt, right=4pt, top=2pt, bottom=2pt}
\newtcolorbox{propbox}[1][]{enhanced, colback=boxProp!5, colframe=boxProp,
  fonttitle=\bfseries\small,
  title={Proposition\ifstrempty{#1}{}{ -- #1}},
  boxrule=0.5pt, arc=2pt, left=4pt, right=4pt, top=2pt, bottom=2pt}
\newtcolorbox{defbox}[1][]{enhanced, colback=boxDef!5, colframe=boxDef,
  fonttitle=\bfseries\small,
  title={Definition\ifstrempty{#1}{}{ -- #1}},
  boxrule=0.5pt, arc=2pt, left=4pt, right=4pt, top=2pt, bottom=2pt}
\newtcolorbox{exbox}[1][]{enhanced, colback=boxEx!5, colframe=boxEx,
  fonttitle=\bfseries\small,
  title={Example\ifstrempty{#1}{}{ -- #1}},
  boxrule=0.5pt, arc=2pt, left=4pt, right=4pt, top=2pt, bottom=2pt}


\title[\moduleNom\ -- Ch.4]{Non-stationary Time Series}
\subtitle{Operators $\cdot$ Unit roots $\cdot$ ADF / KPSS $\cdot$ ARIMA / SARIMA}
\author{\auteurNom \\ \footnotesize\auteurEmail}
\institute{\auteurInstitution}
\date{\anneeUniversitaire}

\begin{document}

\begin{frame}\titlepage\end{frame}

\begin{frame}{Learning objectives}
  By the end of this chapter, you will be able to:
  \begin{itemize}
    \item recognise the limits of the stationarity assumption and
          the danger of \emph{spurious regressions};
    \item distinguish \emph{trend stationary} (TS) from
          \emph{difference stationary} (DS, with a unit root);
    \item manipulate the \emph{backshift} $B$ and \emph{difference}
          $\Delta = 1 - B$ operators;
    \item run the \emph{ADF} and \emph{KPSS} unit-root tests and read
          their conclusions correctly;
    \item formulate \emph{ARIMA(p, d, q)} and
          \emph{SARIMA(p, d, q)(P, D, Q)$_s$} models.
  \end{itemize}
\end{frame}

\begin{frame}{Chapter outline}
  \begin{enumerate}
    \item Why non-stationarity matters (spurious regression)
    \item Backshift and difference operators
    \item Trend-stationary vs difference-stationary
    \item Unit-root tests: ADF, KPSS, Phillips--Perron
    \item ARIMA(p, d, q)
    \item Seasonal ARIMA: SARIMA(p, d, q)(P, D, Q)$_s$
  \end{enumerate}
\end{frame}

% =====================================================================
\section{Why non-stationarity matters}

\begin{frame}{1.~Limits of the stationarity assumption}
  Most real-world series — GDP, asset prices, populations,
  cumulative indices — are \emph{not} stationary:
  \begin{itemize}
    \item their mean drifts upward / downward over time;
    \item their variance often grows;
    \item shocks can have permanent effects.
  \end{itemize}
  \medskip
  \begin{warnbox}[Spurious regression (Granger \& Newbold, 1974)]
    Regressing one random walk on another \emph{independent} random
    walk yields a high $R^2$ and a ``significant'' $t$-statistic,
    even though there is no causal link. This famous fallacy is
    avoided only by working with stationary (or co-integrated)
    series.
  \end{warnbox}
\end{frame}

\begin{frame}{1.~Two flavours of non-stationarity}
  \begin{defbox}[Trend stationary (TS)]
    $X_t = f(t) + \varepsilon_t$ with $\varepsilon_t$ stationary.
    Removing the deterministic trend $f(t)$ recovers a stationary
    process. Shocks have \emph{transient} effects.
  \end{defbox}
  \medskip
  \begin{defbox}[Difference stationary (DS)]
    $\Delta X_t = X_t - X_{t-1}$ is stationary, but $X_t$ has a
    \emph{unit root} (e.g.\ random walk). Shocks are \emph{permanent}.
  \end{defbox}
  \medskip
  \begin{infobox}
    Same plot can be \emph{either} TS or DS — visual inspection alone
    is not enough. Hence the need for a formal \emph{unit-root test}.
  \end{infobox}
\end{frame}

% =====================================================================
\section{Backshift and difference operators}

\begin{frame}{2.~Backshift operator $B$}
  \begin{defbox}[Backshift]
    $B X_t \;=\; X_{t-1}$, $B^k X_t \;=\; X_{t-k}$.
  \end{defbox}
  \medskip
  \textbf{Algebra of $B$.} For any polynomials $\phi(z), \theta(z)$:
  \[
    \phi(B)\, \theta(B) = \theta(B)\, \phi(B), \quad \text{and}
    \quad
    (1 - \phi B) X_t = X_t - \phi X_{t-1}.
  \]
  \medskip
  \textbf{Compact AR(1):} $X_t = \phi X_{t-1} + \varepsilon_t$
  $\Leftrightarrow$ $(1 - \phi B) X_t = \varepsilon_t$.
  \medskip
  \begin{itemize}
    \item Stationarity $\Leftrightarrow$ all roots of $\phi(z) = 0$
          lie \emph{outside} the unit disk.
    \item $\phi(z) = 1 - \phi z$ has a root at $z = 1/\phi$:
          stationary iff $\abs\phi < 1$.
  \end{itemize}
\end{frame}

\begin{frame}{2.~Difference operator $\Delta$}
  \begin{defbox}[Difference]
    $\Delta X_t \;=\; X_t - X_{t-1} \;=\; (1 - B) X_t$.
  \end{defbox}
  \medskip
  \textbf{Higher orders:}
  \[
    \Delta^2 X_t = \Delta(\Delta X_t) = X_t - 2 X_{t-1} + X_{t-2}
                = (1 - B)^2 X_t.
  \]
  \medskip
  \textbf{Seasonal difference:}
  \[
    \Delta_s X_t = X_t - X_{t-s} = (1 - B^s) X_t.
  \]
  \medskip
  \begin{infobox}
    A series is \emph{integrated of order $d$}, written
    $X_t \sim I(d)$, if $\Delta^d X_t$ is stationary but
    $\Delta^{d-1} X_t$ is not.
  \end{infobox}
\end{frame}

\begin{frame}{2.~Random walk: archetypal $I(1)$ process}
  $X_t = X_{t-1} + \varepsilon_t$ with $\varepsilon_t \sim \WN(0, \sigma^2)$.
  \[
    \Var[X_t] = t\,\sigma^2 \to \infty,
    \qquad
    \Delta X_t = X_t - X_{t-1} = \varepsilon_t \in \WN(0, \sigma^2).
  \]
  \medskip
  \textbf{With drift:} $X_t = c + X_{t-1} + \varepsilon_t$
  $\Rightarrow$ $\Delta X_t = c + \varepsilon_t$ — stationary mean $c$.
  \medskip
  \begin{warnbox}
    A random walk \emph{does not} fluctuate around a fixed level. It
    wanders, with a variance that grows linearly in $t$. This is the
    null hypothesis of the ADF test.
  \end{warnbox}
\end{frame}

% =====================================================================
\section{Unit-root tests}

\begin{frame}{3.~Augmented Dickey--Fuller (ADF) test}
  Model: $X_t = \rho X_{t-1} + \varepsilon_t$. Subtracting $X_{t-1}$:
  \[
    \Delta X_t \;=\; \gamma\, X_{t-1} + \varepsilon_t,
    \qquad \gamma = \rho - 1.
  \]
  \emph{Augmented} version (adds lags of $\Delta X_t$ to whiten residuals):
  \[
    \Delta X_t \;=\; \gamma\, X_{t-1} + \sum_{i = 1}^{p} \alpha_i\, \Delta X_{t-i} + \varepsilon_t.
  \]
  \medskip
  \begin{defbox}[Hypotheses]
    $H_0$: $\gamma = 0$ — \textbf{unit root, non-stationary}.\\
    $H_1$: $\gamma < 0$ — stationary (mean-reverting).
  \end{defbox}
  \medskip
  Reject $H_0$ if the $t$-statistic for $\widehat\gamma$ is more
  negative than the Dickey--Fuller critical values
  ($\sim -2.86$ at 5\% with intercept; \emph{not} the standard normal!).
\end{frame}

\begin{frame}{3.~ADF: practical decision rule and traps}
  \begin{itemize}
    \item Use the \emph{augmented} version with $p$ chosen by AIC/BIC.
    \item Three flavours: no constant, with constant, with constant and
          trend. Choose based on the visual pattern of the series.
    \item Power is low when $\rho$ is close to $1$ (near unit root).
  \end{itemize}
  \medskip
  \begin{warnbox}
    Be careful with the \emph{direction of the conclusion}:
    \begin{itemize}
      \item $p$-value $< 0.05$ $\Rightarrow$ \textbf{reject} $H_0$
            $\Rightarrow$ stationary;
      \item $p$-value $> 0.05$ $\Rightarrow$ \emph{do not reject} $H_0$
            $\Rightarrow$ non-stationarity is \emph{not ruled out}
            (this is \emph{not} ``the series is stationary'').
    \end{itemize}
  \end{warnbox}
\end{frame}

\begin{frame}{3.~KPSS test (opposite null hypothesis)}
  \begin{defbox}[KPSS — Kwiatkowski et al., 1992]
    $H_0$: \textbf{stationarity} (around a level or a trend).\\
    $H_1$: unit-root non-stationarity.
  \end{defbox}
  \medskip
  Reject $H_0$ at $5\%$ if the test statistic exceeds $0.463$
  (level) or $0.146$ (trend) — depending on the version.
  \medskip
  \begin{infobox}[Practical use: ADF + KPSS together]
    \begin{itemize}
      \item ADF rejects + KPSS does not reject $\Rightarrow$ stationary.
      \item ADF does not reject + KPSS rejects $\Rightarrow$ non-stationary.
      \item ADF rejects + KPSS rejects $\Rightarrow$ inconclusive
            (probably structural break).
      \item ADF does not reject + KPSS does not reject $\Rightarrow$
            insufficient information (small sample).
    \end{itemize}
  \end{infobox}
\end{frame}

\begin{frame}[fragile]{3.~ADF and KPSS in Python}
\begin{lstlisting}
from statsmodels.tsa.stattools import adfuller, kpss

# Augmented Dickey-Fuller
adf_stat, adf_p, adf_lag, adf_n, adf_crit, _ = adfuller(y)
print(f"ADF stat={adf_stat:.3f}, p-value={adf_p:.4f}")

# KPSS (level by default)
kpss_stat, kpss_p, kpss_lag, kpss_crit = kpss(y, regression="c")
print(f"KPSS stat={kpss_stat:.3f}, p-value={kpss_p:.4f}")
\end{lstlisting}
\end{frame}

% =====================================================================
\section{ARIMA}

\begin{frame}{4.~ARIMA(p, d, q)}
  \begin{defbox}[ARIMA(p, d, q)]
    A process $X_t$ is \emph{ARIMA(p, d, q)} if its $d$-th difference
    is a stationary ARMA(p, q) :
    \[
      \phi(B)\, (1 - B)^d\, X_t \;=\; \theta(B)\, \varepsilon_t,
      \qquad \varepsilon_t \sim \WN(0, \sigma^2),
    \]
    with $\phi(z) = 1 - \phi_1 z - \dots - \phi_p z^p$ and
    $\theta(z) = 1 + \theta_1 z + \dots + \theta_q z^q$.
  \end{defbox}
  \medskip
  \begin{itemize}
    \item Bring $X_t$ to stationarity by differencing $d$ times.
    \item Fit ARMA(p, q) on the differenced series.
    \item Forecasts on the original scale are obtained by integrating
          back ($X_t = X_{t-1} + \Delta X_t$).
  \end{itemize}
\end{frame}

\begin{frame}{4.~Box--Jenkins methodology}
  \begin{enumerate}
    \item \textbf{Plot} $X_t$. Identify trend, seasonality, level shifts.
    \item \textbf{Test} for unit roots (ADF + KPSS). If non-stationary,
          difference until stationary $\Rightarrow$ choose $d$.
    \item \textbf{Identify} $(p, q)$ from the ACF/PACF of the
          differenced series (Box--Jenkins signature, Chapter 3).
    \item \textbf{Estimate} the ARMA(p, q) coefficients (MLE).
    \item \textbf{Diagnose}: Ljung--Box on residuals, plot of residuals.
    \item \textbf{Forecast} and integrate back.
  \end{enumerate}
\end{frame}

\begin{frame}{4.~Example: ARIMA(1, 1, 1)}
  Definition: $\phi(B)(1 - B) X_t = \theta(B) \varepsilon_t$ with
  $\phi(B) = 1 - \phi B$ and $\theta(B) = 1 + \theta B$.
  \medskip
  Expanded:
  \[
    (1 - \phi B)(1 - B) X_t = (1 + \theta B) \varepsilon_t.
  \]
  Letting $W_t = \Delta X_t$:
  \[
    W_t = \phi W_{t-1} + \varepsilon_t + \theta \varepsilon_{t-1}.
  \]
  $W_t$ is an ARMA(1,1); $X_t$ is its cumulative sum starting from
  $X_0$.
\end{frame}

% =====================================================================
\section{Seasonal ARIMA (SARIMA)}

\begin{frame}{5.~SARIMA(p, d, q)(P, D, Q)$_s$}
  \begin{defbox}[SARIMA]
    \[
      \phi(B)\, \Phi(B^s)\,
      (1 - B)^d\, (1 - B^s)^D\, X_t
      \;=\; \theta(B)\, \Theta(B^s)\, \varepsilon_t,
    \]
    with seasonal AR/MA polynomials $\Phi, \Theta$ in $B^s$.
  \end{defbox}
  \medskip
  \textbf{Six orders}:
  \begin{itemize}
    \item \textbf{Non-seasonal}: $p$ (AR), $d$ (regular differencing), $q$ (MA);
    \item \textbf{Seasonal}: $P$ (SAR), $D$ (seasonal differencing), $Q$ (SMA);
    \item \textbf{Period}: $s$ (e.g.\ $12$ for monthly data).
  \end{itemize}
\end{frame}

\begin{frame}{5.~Example: SARIMA(1, 1, 1)(1, 1, 1)$_{12}$}
  \[
    (1 - \phi B)(1 - \Phi B^{12})(1 - B)(1 - B^{12}) X_t
    = (1 + \theta B)(1 + \Theta B^{12}) \varepsilon_t.
  \]
  \medskip
  \begin{itemize}
    \item Differences once at lag 1 \emph{and} once at lag 12.
    \item Captures both short-term persistence and yearly seasonality.
    \item Standard for monthly economic, climate and sales data.
  \end{itemize}
\end{frame}

\begin{frame}[fragile]{5.~SARIMA in Python}
\begin{lstlisting}
from statsmodels.tsa.statespace.sarimax import SARIMAX

model = SARIMAX(y,
                order=(1, 1, 1),
                seasonal_order=(1, 1, 1, 12))
fit = model.fit(disp=False)
print(fit.summary())

# 12-month ahead forecast
forecast = fit.get_forecast(steps=12)
mean = forecast.predicted_mean
ci = forecast.conf_int(alpha=0.05)
\end{lstlisting}
\end{frame}

\begin{frame}{Up next}
  \textbf{Chapter 5.} ARMA models — estimation, diagnostics,
  forecasting in detail.
  \medskip
  \begin{itemize}
    \item AR(p), MA(q), ARMA(p, q) algebra.
    \item Yule--Walker, MLE estimation.
    \item Order selection (AIC, BIC).
    \item Forecasting and prediction intervals.
  \end{itemize}
  \medskip
  \begin{center}\Large\bfseries Thank you.\end{center}
  \begin{center}\footnotesize\auteurNom\ -- \auteurEmail\ -- \auteurInstitution\end{center}
\end{frame}

\end{document}
